\input{/home/mbenjamin/Documents/Lycee/LateX/head.tex}

% Mise en page
\usepackage[left=2.0cm, right=2.0cm, top=1.5cm, bottom=1.5cm]{geometry}

% Numerotation des pages / la derniere
\makeatletter
\renewcommand{\@evenfoot}%
        {\thepage \slash \pageref{LastPage}}
\renewcommand{\@oddfoot}{\@evenfoot}
\makeatother

% En-tête
\renewcommand{\chapitre}{Statistique}
\renewcommand{\seance}{Statistiques à 2 variables et coefficient de correlation}
\renewcommand{\classe}{CIEL 2}



\begin{document}

\begin{titre}
Statistique descriptive \\

Série statistique à deux variables
\end{titre}


On donne ci-dessous le nombre d'étudiants diplômés d'un Master en mathématiques et statistiques aux États-Unis$^1$ ainsi que le nombre de programmeurs informatiques en Californie$^2$ entre 2012 et 2021.

\smallskip

\hfill {\footnotesize $^1$ étude du Centre National pour les Statistiques sur l'Éducation (NCES) des États-Unis.}

\hfill {\footnotesize $^2$ étude du Bureau des Statiques sur l'Emploi (BLS) des États-Unis.}

\begin{center}
\begin{tblr}{
    colspec={lXXXXXXXXXX},
    stretch=1,
    cells={valign=m},
    vlines,hlines,
    row{1} = {bg = darkWhite},
    column{1} = {bg = darkWhite}
}
	Année & 2012 & 2013 & 2014 & 2015 & 2016 & 2017 & 2018 & 2019 & 2020 & 2021 \\
	Diplômés & 6246 & 6957 & 7273 & 7589 & 8451 & 9082 & 10443 & 11382 & 12041 & 12583 \\
	Programmeurs & 41540 & 38750 & 37990 & 38650 & 	37570 & 34050 & 29740 & 24400 & 21800 & 23960
\end{tblr}
\end{center}

On souhaite évaluer si il existe un lien (une corrélation) entre ces données.

\begin{definition}
\underline{\bf Définition :}

\smallskip

On considère deux caractéristiques $x$ et $y$ d'une même population.

On définit une \textbf{série statistique à deux variables} $x$ et $y$ prenant les valeurs $x_1$, ..., $x_i$, ... $x_n$

et $y_1$, ..., $y_i$, ... $y_n$ par un tableau comme celui-ci :

\medskip

\begin{tblr}{
    colspec={lXXXXX},
    stretch=1,
    cells={valign=m, halign=c},
    vlines,hlines,
    column{1} = {bg = darkGreen!25!white}
}
	Première variable $x$ & $x_1$ & \dots & $x_i$ & \dots & $x_n$ \\
	Deuxième variable $y$ & $y_1$ & \dots & $y_i$ & \dots & $y_n$
\end{tblr}
\end{definition}

\vspace*{-0.5cm}

{\color{NavyBlue}\section{Nuage de points}}


\begin{minipage}{0.65\lw}
\begin{definition}
\underline{\bf Définition :}

\smallskip

On parle d'un nuage de point représentant une série statistique à deux variables lorsqu'on place, dans un repère orthogonal, l'ensemble des points de coordonnées $(x_i~;~y_i)$, c'est-à-dire tous les points $(x_n~;~y_n)$, $(x_n~;~y_n)$ jusqu'à $(x_n~;~y_n)$.
\end{definition}

{\it Sur Géogebra :}

\end{minipage}
\hfill
\begin{minipage}{0.3\lw}
	\includegraphics[width=\lw]{nuage0}
	
	\bigskip
\end{minipage}

\texttt{Affichage > Tableur}, remplir les données,

puis \texttt{Statistiques à deux variables > Analyse}\par  \medskip

\begin{definition}
\underline{\bf Définition :} Point moyen

\smallskip

On note $\overline{x}$ et $\overline{y}$ les moyennes respectives de la série $x$ et de la série $y$.

Le point moyen de la série statistique à deux variables est alors le point de coordonnées $(\overline{x}~;~\overline{y})$.
\end{definition}


\begin{minipage}{0.65\lw}
\smallskip

Exemple :

\scriptsize

$$\overline{x}=\dfrac{6246+6957+7273+7589+8451+9082+10443+11382+12041+12583}{10} $$

$$\overline{y}=\dfrac{41540+38750+37990+38650+37570+34050+29740+24400+21800+23960}{10}$$
\end{minipage}
\hfill
\begin{minipage}{0.3\lw}
	\includegraphics[width=\lw]{nuage0pm}
\end{minipage}

\medskip

{\color{NavyBlue}\section{Ajustement affine}}

\subsection{Droite de regression}

Lorsque les points d’un nuage sont sensiblement alignés, on peut
construire une droite, appelé droite d’ajustement, ou \textbf{droite de régression},
passant « au plus près » de ces points.

\medskip

On pourra toujours choisir plusieurs droite, l'objectif sera alors de minimiser les écarts avec les points. Plusieurs techniques sont possibles.

\begin{center}
	\hfill \includegraphics[width=0.35\lw]{nuage0droite}
	\hfill \includegraphics[width=0.35\lw]{nuage0droite2}
	\hfill ~
\end{center}

\subsection{Interpolation / Extrapolation}

Pour un nuage de points donné, on peut (à l'aide d'un ajustement affine par exemple) déterminer des valeurs << manquantes >>, qui seraient à l'intérieur des données étudiées : on parle d'\textbf{interpolation}.

\medskip

On peut aussi << prolonger >> le nuage en envisageant des valeurs qui seraient en dehors des données étudiées : on parle alors d'\textbf{extrapolation}.

\subsection{Methode de Mayer}

On souhaite approcher un nuage de points par une droite, ce qui nécessite deux points.

\bigskip

\textbf{La droite de Mayer} est une approximation du nuage de points qui passe par deux points moyens :
\begin{itemize}[$\hookrightarrow$]
	\item On coupe la série statistique en deux sous séries d'effectifs similaires.
	\item On calcule les coordonnées du point moyen PM$_1$ de la première sous-série obtenue.
	\item On calcule les coordonnées du point moyen PM$_2$ de la deuxième sous-série obtenue.
	\item La droite de Mayer est la droite (PM$_1$ PM$_2$)
\end{itemize}

\medskip

\begin{center}
	\includegraphics[width=0.5\lw]{nuage0mayer}
\end{center}

\newpage

\entete

\vspace*{-0.6cm}

\subsection{Methode des moindres carrés}

\begin{definition}
\underline{\bf Définition :} variance

\smallskip

La variance d'une série statistique $x$, notée $var(x)$ est la moyenne des écarts à la moyenne au carré :\par \medskip
\hfill$\displaystyle var(x)=\dfrac{(x_1-\overline{x})^2+(x_2-\overline{x})^2+\dots+(x_i-\overline{x})^2+\dots+(x_n-\overline{x})^2}{n}=\dfrac{1}{n}\sum_{i=1}^{n}(x_i-\overline{x})^2$\hfill~
\end{definition}


Lorsqu'on trace une regression liénaire avec la méthode des moindres carrés, on cherche à minimiser la somme des carrés des distances entre les points et la droite qu'on va tracer ; le plus souvent pour $y$, mais on peut aussi le faire pour $x$.

\medskip

\hfill
\includegraphics[width=0.3\lw]{ex_decouverte_droite_regression2}
\hfill
\includegraphics[width=0.3\lw]{ex_decouverte_droite_regression2b}
\hfill
\includegraphics[width=0.3\lw]{regression_yx}
\hfill ~

\vspace*{0.5cm}

La droite de régression par la méthode des moindres carrés passe par le point moyen de la série statistique à deux variables.

\begin{definition}
\underline{\bf Définition :} covariance

\smallskip

On appelle covariance de la série statistique à deux variables, notée $cov(x,y)$, le nombre :
$$\displaystyle \dfrac{(x_1-\overline{x})(y_1-\overline{y})+\dots+(x_i-\overline{x})(y_i-\overline{y})+\dots+(x_n-\overline{x})(y_n-\overline{y})}{n}=\dfrac{1}{n}\sum_{i=1}^{n}(x_i-\overline{x})(y_i-\overline{y})$$
\end{definition}

\vrule
\hfill
\begin{minipage}{0.985\lw}
	\it Remarque : En remplaçant $y$ par $x$ dans la formule de la covariance, on voit bien l'analogie entre la variance d'une variable $x$ et la covariance de deux variables $x$ et $y$.
\end{minipage}

\begin{propriete}
{\bf Propriété :} 

$\displaystyle cov(x,y)=\dfrac{(x_1\times y_1)+\dots+(x_i\times y_i)+...+(x_n\times y_n)}{n}-\overline{x}\times\overline{y}=\dfrac{1}{n}\sum_{i=1}^{n}x_i\times y_i-\overline{x}\times\overline{y}$
\end{propriete}


Soit une série statistique à deux variables $x$ et $y$ de covariance $cov(x,y)$ et dont la variable $x$ a pour moyenne $\overline{x}$ et pour variance $var(x)$, et dont la variable $y$ a pour moyenne $\overline{y}$ et pour variance $var(y)$ :

\begin{itemize}[$\bullet$]
	\item La droite de régression de $y$ en $x$ possède une équation réduite de la forme $y=ax+b$ où :
	\begin{itemize}[$\hookrightarrow$]
		\item le coefficient directeur $a=\dfrac{cov(x,y)}{var(x)}$
		\item $b$ vérifie l'équation $\overline{y}=a\overline{x}+b$ (puisque cette droite
                de régression passe par le point moyen), soit $b=\overline{y}-a\overline{x}$
	\end{itemize}		
    \item La droite de régression de $x$ en $y$ possède une équation réduite de la forme $x=a'y+b'$ où :
	\begin{itemize}[$\hookrightarrow$]
		\item le coefficient directeur $a'=\dfrac{cov(x,y)}{var(y)}$
		\item $b'=\overline{x}-a'\overline{y}$
	\end{itemize}	    
\end{itemize}

\newpage



{\color{NavyBlue}\section{Coefficient de correlation}}

\begin{definition}
\underline{\bf Définition :} coefficient de correlation

\smallskip

Le coefficient de corrélation linéaire d'une série statistique à deux variables $x$ et $y$, noté $r$, est le nombre nombre défini par $r=\dfrac{cov(x,y)}{\sqrt{var(x)\times var(y)}}$
\end{definition}

\underline{\bf Interprétation :}\par \medskip

Le coefficient de corrélation $r$ est un nombre compris entre $-1$ et $1$ qui mesure la
relation entre les deux variables $x$ et $y$.

\medskip

Plus le coefficient est proche des valeurs extrêmes $-1$ et $1$, plus la corrélation linéaire entre les variables est forte.

\begin{itemize}[$\hookrightarrow$]
	\item Si $r > 0$, les valeurs prises par $y$ ont tendance à croître quand les valeurs de $x$ augmentent,
	\item si $r < 0$, les valeurs prises par $y$ ont tendance à décroître quand les valeurs de $x$ augmentent,
	\item si $r = 0$, les variations des variables $x$ et $y$ sont indépendantes.
\end{itemize}

\bigskip

\underline{Exemples de nuages de points en fonction du coefficient de correlation :}

\begin{center}
	\includegraphics[width=0.8\lw]{coeffs}
\end{center}

\begin{tcolorbox}
{\bf Correlation vs lien de cause à effet :} {\it cum hoc sed non propter hoc}

\medskip

{\it ATTENTION}, des données fortement corrélées ne signifie que les données ont un lien de cause à effet. Certaines correlation peuvent résulter du hasard et n'avoir aucun lien.

\bigskip

Il peut d'ailleurs être amusant de volontairement choisir des données totalement indépendantes mais qui ont une correlation forte pour inventer des liens :

\smallskip

le nombre de nids de cigognes et celui des naissances humaines par exemple !
\end{tcolorbox}

\end{document}


